Inserendo un campione di germanio drogato $ p $ all'interno di un elettromagnete studiamo l'effetto Hall: la formazione di una differenza di potenziale tra i lati della sonda.  \\
Valutando gli andamenti della tensione di Hall $ V_H $ contro la corrente $ i $ e il campo magnetico $ B $, diamo una stima del coefficiente di Hall $ R_H = (0.0079 \pm 0.0006) \unit{m^3 / C} $ del campione, e della concentrazione dei portatori di carica maggioritari $ \rho = (127 \pm 9) \unit{C / m^3} $.
Infine, studiando anche la caratteristica $ I(V) $ della sonda siamo in grado di fornire una stima della mobilità dei portatori al suo interno $ \mu = (0.24 \pm 0.03) \unit{m^2 / (V \cdot s)} $. 
